\section{Multilinear Algebra}

\noindent\textbf{Exercise~\ref{ex:multilinear-1}.}\label{sol:multilinear-1}

We need to check it meets the 8 axioms.
\benr
    \item Commutative w.r.t. $\oplus$:
    \bse
        (a,b,c) \oplus (d,e,f) = (a+d,b+e,c+f) = (d+a,b+e,f+c) = (d,e,f) \oplus (a,b,c).
    \ese
    \item Associative w.r.t. $\oplus$:
    \bse
        \begin{split}
            \big[ (a,b,c)\oplus(d,e,f) \big] \oplus (h,i,j) & = (a+d,b+e,c+f) \oplus (h,i,j) \\
            & = (a+d+h,b+e+i,c+f+j) \\
            & = (a,b,c) \oplus (d+h,e+i,f+j) \\
            & = (a,b,c)\oplus \big[ (d,e,f)\oplus(h,i,j) \big].
        \end{split}
    \ese
    \item Neutral element w.r.t. $\oplus$: this is clearly just $(0,0,0)$.
    \item Inverse w.r.t. $\oplus$: this is clearly just $(-a,-b,-c)$. Note for this to be true it is important that we take the whole real line, not just the positive reals (otherwise $(-a,-b,-c)\notin V$.
    \item Associative w.r.t. $\odot$:
    \bse
        \begin{split}
            \lambda \odot \big[\mu\odot(a,b,c)\big] & = \lambda \odot (\mu\cdot a,\mu\cdot b, \mu\cdot c) \\
            & = (\lambda \odot\mu\cdot a,\lambda \odot\mu\cdot b, \lambda \odot\mu\cdot c) \\
            & = \big[\lambda \odot \mu] \odot (a,b,c).
        \end{split}
    \ese
    \item Distributive 1:
    \bse
        \begin{split}
            (\lambda + \mu)\odot (a,b,c) & = \big( (\lambda + \mu)\cdot a, (\lambda + \mu)\cdot b,(\lambda + \mu)\cdot c\big) \\
            & = (\lambda \cdot a, \lambda \cdot b, \lambda \cdot c) \oplus (\mu\cdot a,\mu\cdot b, \mu \cdot c) \\
            & = \lambda \odot (a,b,c) \oplus \mu \odot (a,b,c).
        \end{split}
    \ese
    \item Distributive 2:
    \bse
        \begin{split}
            \lambda \odot \big[(a,b,c)\oplus(d,e,f)\big] & = \lambda \odot (a+d,b+e,c+f) \\
            & = \big( \lambda\cdot(a+d),\lambda\cdot(b+e), \lambda\cdot(c+f)\big) \\
            & = (\lambda\cdot a, \lambda\cdot b, \lambda\cdot c) + (\lambda\cdot d, \lambda\cdot e, \lambda\cdot f) \\
            & = \lambda \odot (a,b,c) \oplus \lambda \odot (d,e,f).
        \end{split}
    \ese
    \item Unitary w.r.t. $\odot$: Clearly $1\odot(a,b,c)=(a,b,c)$.
\een

\bigskip
\noindent\textbf{Exercise~\ref{ex:multilinear-2}.}\label{sol:multilinear-2}

Consider
\bse
    \begin{split}
        d\big( \lambda\odot (a,b,c) \oplus (d,e,f)\big) & = d\big( ([\lambda \cdot a]+d, [\lambda \cdot b]+e,[\lambda \cdot c]+f)\big) \\
        & = ([\lambda \cdot b]+e, 2[\lambda \cdot c]+2f, 0) \\
        & = \lambda\odot(b, 2c,0) \oplus (e,2f) \\
        & = \lambda\odot d\big((a,b,c)\big) \oplus d\big((d,e,f)\big),
    \end{split}
    so yes it's linear.
\ese

\bigskip
\noindent\textbf{Exercise~\ref{ex:multilinear-3}.}\label{sol:multilinear-3}

This follows immediately from \Cref{thrm:CompositionOfLinearMaps}, however if you want you can show it explicitly (we won't here to save me writing.)

\bigskip
\noindent\textbf{Exercise~\ref{ex:multilinear-4}.}\label{sol:multilinear-4}

The linearity calculation is exactly the same method as above, so to save me typing I'm just going to say the answer: yes it is. Being a linear map from $V$ to $\R$ it is, by definition, an element of $V^*$.

\bigskip
\noindent\textbf{Exercise~\ref{ex:multilinear-6}.}\label{sol:multilinear-6}

Recall that $\del:P\to P$ is defined as the derivative, i.e. $\del(p) = p'$. Clearly for this question we want to focus on the map $\del : P_2\to P_2$. Recall the definition
\bse
    P_2 := \{ p :\R \to \R \, | \, p(x) = a + bx + cx^2\},
\ese
and so we have
\bse
    p'(x) = b + 2cx + 0x^2.
\ese
We instantly see this resembles the map $d: V \to V$ which acts as $d\big((a,b,c)\big) := (b,2c,0)$. We just need to construct the map $j:P_2\to \R^3$. With a bit of thought it is clear that the answer is simply
\bse
    j(p) := (a,b,c),
\ese
where $p = a + bx +cx^2$. The inverse map $j^{-1}:\R^3\to P_2$ is then simply given by
\bse
    \big(j^{-1}(a,b,c)\big)(x) := a+bx+cx^2.
\ese
Direct substitution then gives
\bse
    \big((j\circ \del \circ j^{-1}\big)(a,b,c) := (b,2c,0) =: d\big((a,b,c)\big).
\ese

\bigskip
\noindent\textbf{Exercise~\ref{ex:multilinear-7}.}\label{sol:multilinear-7}

We simply uses the duality of the basis elements, namely $e_i(\epsilon^j) = \del^j_i = \epsilon^j(e_i)$. So we have
\bse
    A^{ab} = A(\epsilon^a,\epsilon^b), \qand B_{ab} = B(e_a,e_b).
\ese

\bigskip
\noindent\textbf{Exercise~\ref{ex:multilinear-8}.}\label{sol:multilinear-8}

From direct calculation we have
\bse
    A^{[ab]} := \frac{1}{2}\big(A^{ab}-A^{ba}\big) = - \frac{1}{2} \big( A^{ba} - A^{ab}\big) =: - A^{[ba]}.
\ese
We also have
\bse
    \begin{split}
        A^{[ab]}B_{ab} & := \frac{1}{2}\big(A^{ab}-A^{ba}\big)B_{ab} \\
        & = \frac{1}{2} \big( A^{ab}B_{ab} - A^{ba}B_{ab}\big) \\
        & = \frac{1}{2} \big( A^{ab}B_{ab} - A^{ab}B_{ba}\big) \\
        & = A^{ab}\frac{1}{2}\big(B_{ab}-B_{ba}\big) \\
        & =: A^{ab}B_{[ab]},
    \end{split}
\ese
where we have used the Einstein summation convention to relabel the dummy indices of the second term to get to the third line.

\bigskip
\noindent\textbf{Exercise~\ref{ex:multilinear-9}.}\label{sol:multilinear-9}

Follows exactly like the previous solution.

\bigskip
\noindent\textbf{Exercise~\ref{ex:multilinear-10}.}\label{sol:multilinear-10}

We have
\bse
    \begin{split}
        A^{[ab]}B_{(ab)} & := \frac{1}{2}\big( A^{ab}B_{(ab)} - A^{ba}B_{(ab)} \big) \\
        & = \frac{1}{2}\big( A^{(ab)}B_{ab} - A^{(ba)}B_{ab} \big) \\
        & = \frac{1}{2}\big( A^{(ab)}B_{ab} - A^{(ab)}B_{ab}\big) \\
        & = 0.
    \end{split}
\ese

\bigskip
\noindent\textbf{Exercise~\ref{ex:multilinear-11}.}\label{sol:multilinear-11}

We know a $(1,1)$-tensor is a linear map $T_{\phi}: V^*\times V \lmap \R$. We also know that $\psi\in V^*$ is a map $\psi: V\lmap \R$. So it's clear that we just define
\bse
    T_{\phi}(\sig, v) := \big(\phi(\sig)\big)(v),
\ese
where $\sig\in V^*$ and $v\in V$.

\bigskip
\noindent\textbf{Exercise~\ref{ex:multilinear-12}.}\label{sol:multilinear-12}

This is just the previous exercise in reverse. That is we define
\bse
    \big(\phi_T(\sig)\big)(v) := T(\sig,v),
\ese
which we write as
\bse
    \phi_T(\sig) := T(\sig, \cdot).
\ese

\bigskip
\noindent\textbf{Exercise~\ref{ex:multilinear-13}.}\label{sol:multilinear-13}

\ben[label=(\alph*)]
    \item We have
    \bse
        T_{\phi_T}(\sig,v) := \big(\phi_T(\sig)\big)(v) =: T(\sig,v).
    \ese
    \item We have
    \bse
        \big(\phi_{T_{\phi}}(\sig)\big)(v) := T_{\phi}(\sig,v) =: \big(\phi(\sig)\big(v),
    \ese
    and so
    $\phi_{T_{\phi}} = \phi$.
\een

\bigskip
\noindent\textbf{Exercise~\ref{ex:multilinear-14}.}\label{sol:multilinear-14}

The above exercises have just shown us that there is a unique link between $T$ and $\phi$ and so we can think of them as being isomorphic as maps and as such can be viewed as the same object.

\bigskip
\noindent\textbf{Exercise~\ref{ex:multilinear-15}.}\label{sol:multilinear-15}

First we show $\mathbf{0}$ is linear. For $v,w\in V$ and $\lambda \in \R$,
\bse
    \mathbf{0}(v+w) = 0_W = 0_W + 0_W = \mathbf{0}(v) + \mathbf{0}(w),
\ese
and
\bse
    \mathbf{0}(\lambda \cdot v) = 0_W = \lambda \cdot 0_W = \lambda \cdot \mathbf{0}(v),
\ese
where the second equality uses the fact that $\lambda \cdot 0_W = 0_W$ in any vector space. So $\mathbf{0}\in\Hom(V,W)$.

To show it is the neutral element, let $\varphi \in \Hom(V,W)$. Then for all $v\in V$,
\bse
    (\varphi \oplus \mathbf{0})(v) = \varphi(v) + \mathbf{0}(v) = \varphi(v) + 0_W = \varphi(v),
\ese
and so $\varphi \oplus \mathbf{0} = \varphi$.

\bigskip
\noindent\textbf{Exercise~\ref{ex:multilinear-16}.}\label{sol:multilinear-16}

First we show $\widetilde{\varphi}$ is linear. For $v,w\in V$ and $\lambda\in\R$,
\bse
    \begin{split}
        \widetilde{\varphi}(v+w) & = -\varphi(v+w) \\
        & = -\big(\varphi(v) + \varphi(w)\big) \\
        & = -\varphi(v) + \big(-\varphi(w)\big) \\
        & = \widetilde{\varphi}(v) + \widetilde{\varphi}(w),
    \end{split}
\ese
where the second line uses the linearity of $\varphi$, and
\bse
    \widetilde{\varphi}(\lambda \cdot v) = -\varphi(\lambda\cdot v) = -\big(\lambda \cdot \varphi(v)\big) = \lambda \cdot \big(-\varphi(v)\big) = \lambda \cdot \widetilde{\varphi}(v).
\ese
So $\widetilde{\varphi}\in\Hom(V,W)$.

To show it is the inverse, for all $v\in V$,
\bse
    (\varphi \oplus \widetilde{\varphi})(v) = \varphi(v) + \widetilde{\varphi}(v) = \varphi(v) + \big(-\varphi(v)\big) = 0_W = \mathbf{0}(v),
\ese
and so $\varphi \oplus \widetilde{\varphi} = \mathbf{0}$.
